% Part: methods
% Chapter: proofs
% Section: reading-proofs

\documentclass[../../../include/open-logic-section]{subfiles}

\begin{document}

\olfileid{mth}{prf}{rea}
\olsection{خواندن اثبات‌ها}

اثبات‌هایی که در کتاب‌ها و مقاله‌ها می‌یابید به‌ندرت همهٔ جزئیاتی را
که تا اینجا در مثال‌ها آورده‌ایم بیان می‌کنند. نویسندگان اغلب توجه
خواننده را به اینکه چه هنگام حالت‌ها را جدا می‌کنند، چه هنگام اثباتی
غیرمستقیم می‌آورند، یا اینکه تعریفی را به کار می‌برند، جلب نمی‌کنند.
بنابراین هنگام خواندن اثباتی در یک کتاب، برای فهم آن اغلب باید خودتان
این جزئیات را کامل کنید. این کار همچنین تمرین خوبی برای تسلط بر
حرکت‌های گوناگونی است که در اثبات باید انجام دهید. مثالی را ببینیم.

\begin{prop}[جذب]
برای همهٔ مجموعه‌های $A$ و $B$،
\[
A \cap (A \cup B) = A
\]
\end{prop}

\begin{proof}
اگر $z \in A \cap (A \cup B)$، آنگاه $z \in A$؛ پس
$A \cap (A \cup B) \subseteq A$. اکنون فرض کنید $z \in A$. آنگاه
$z \in A \cup B$ نیز هست و بنابراین $z \in A \cap (A \cup B)$ نیز هست.
\end{proof}

اثبات پیشین قانون جذب بسیار فشرده است. در آن از هیچ‌یک از تعریف‌های
به‌کاررفته نامی برده نمی‌شود، پیش از اثبات چیزی نوشته نمی‌شود «باید
نشان دهیم که»، و مانند آن. بیایید آن را باز کنیم. حکم اثبات‌شده گزاره‌ای
کلی دربارهٔ هر دو مجموعهٔ $A$ و $B$ است، و وقتی اثبات از $A$ یا $B$
نام می‌برد، آن‌ها متغیرهایی برای مجموعه‌های دلخواه‌اند. گزارهٔ کلی که
اثبات برقرار می‌کند همان چیزی است که برای اثبات برابری مجموعه‌ها لازم
است: هر !!{element} از سمت چپ تساوی، !!a{element} از سمت راست است و
برعکس.

\begin{quote}
«اگر $z \in A \cap (A \cup B)$، آنگاه $z \in A$؛ پس
  $A \cap (A \cup B) \subseteq A$.»
\end{quote}

این نیمهٔ نخست اثبات تساوی است: نشان می‌دهد اگر~$z$ دلخواه
!!a{element} از سمت چپ باشد، !!a{element} از سمت راست نیز هست؛ یعنی
$A \cap (A \cup B) \subseteq A$. فرض کنید
$z \in A \cap (A \cup B)$. چون $z$، !!{element} از اشتراک دو مجموعه
است اگر و تنها اگر !!{element} از هر دو مجموعه باشد، می‌توانیم نتیجه
بگیریم $z \in A$ و نیز $z \in A \cup B$. به‌ویژه $z \in A$، و همین
چیزی بود که می‌خواستیم نشان دهیم. چون برای نیمهٔ نخست فقط همین کار لازم
بود، می‌دانیم باقی اثبات باید نیمهٔ دوم را اثبات کند؛ یعنی نشان دهد
$A \subseteq A \cap (A \cup B)$.

\begin{quote}
«اکنون فرض کنید $z \in A$. آنگاه $z \in A \cup B$ نیز هست و بنابراین
  $z \in A \cap (A \cup B)$ نیز هست.»
\end{quote}

با فرض $z \in A$ آغاز می‌کنیم، زیرا نشان می‌دهیم برای هر~$z$، اگر
$z \in A$ آنگاه $z \in A \cap (A \cup B)$. برای نشان‌دادن
$z \in A \cap (A \cup B)$، باید بنا بر تعریف «$\cap$» نشان دهیم
(۱)~$z \in A$ و نیز (۲)~$z \in A \cup B$. اینجا (۱) همان فرض ماست؛
پس چیزی بیشتر برای اثبات وجود ندارد و به همین دلیل اثبات دوباره از آن
نام نمی‌برد. برای (۲)، به یاد آورید که $z$، !!a{element} از اجتماع
مجموعه‌هاست اگر و تنها اگر !!{element} از دست‌کم یکی از آن مجموعه‌ها
باشد. چون $z \in A$ و $A \cup B$ اجتماع $A$ و $B$ است، این شرط اینجا
برقرار است. پس $z \in A \cup B$. هم (۱)~$z \in A$ و هم
(۲)~$z \in A \cup B$ را نشان داده‌ایم؛ بنابراین بنا بر تعریف «$\cap$»،
$z \in A \cap (A \cup B)$. اثبات از این تعریف‌ها نام نمی‌برد؛ فرض
می‌شود خواننده پیش‌تر آن‌ها را درونی کرده است. اگر چنین نکرده‌اید، باید
بازگردید و آن‌ها را به خود یادآوری کنید. سپس باید تشخیص دهید چرا از
$z \in A$، $z \in A \cup B$ نتیجه می‌شود، و چرا از $z \in A$ و
$z \in A \cup B$، $z \in A \cap (A \cup B)$ نتیجه می‌شود.

این هم نسخه‌ای دیگر از اثبات بالا که در آن همه‌چیز صریح شده است:
\begin{proof}{}
[بنا بر تعریف $=$ برای مجموعه‌ها، برای اثبات
  $A \cap (A \cup B) = A$ باید نشان دهیم (الف)
  $A \cap (A \cup B) \subseteq A$ و (ب)
  $A \subseteq A \cap (A \cup B)$. (الف): بنا بر تعریف $\subseteq$،
  باید نشان دهیم اگر $z \in A \cap (A \cup B)$، آنگاه $z \in A$.]
اگر $z \in A \cap (A \cup B)$، آنگاه $z \in A$
[زیرا بنا بر تعریف $\cap$، $z \in A \cap (A \cup B)$ اگر و تنها اگر
  $z \in A$ و $z \in A \cup B$]؛ پس
$A \cap (A \cup B) \subseteq A$. [(ب): بنا بر تعریف $\subseteq$ باید
  نشان دهیم اگر $z \in A$، آنگاه $z \in A \cap (A \cup B)$.]
اکنون فرض کنید [(۱)] $z \in A$. آنگاه [(۲)] $z \in A \cup B$ نیز هست
[زیرا بنا بر (۱)، $z \in A$ یا $z \in B$، که بنا بر تعریف $\cup$ یعنی
  $z \in A \cup B$]، و بنابراین $z \in A \cap (A \cup B)$ نیز هست
[زیرا تعریف $\cap$ مستلزم $z \in A$، یعنی (۱)، و $z \in A \cup B$،
  یعنی (۲)، است].
\end{proof}

\begin{prob}
اثبات زیر از $A \cup (A \cap B) = A$ را بسط دهید: همهٔ الگوهای
استنتاجی به‌کاررفته را نام ببرید، توضیح دهید چرا هر گام از فرض‌ها یا
ادعاهای پیش‌تر اثبات‌شده نتیجه می‌شود، و مشخص کنید در کجا باید به کدام
تعریف‌ها متوسل شویم.
\begin{proof}
  اگر $z \in A \cup (A \cap B)$، آنگاه $z \in A$ یا
  $z \in A \cap B$. اگر $z \in A \cap B$، آنگاه $z \in A$. هر
  $z \in A$ در $A \cup (A \cap B)$ نیز هست.
\end{proof}
\end{prob}

\end{document}
